\labsection{8}{BSTs, heaps, and balancing structures}{sec:lab08}

\begin{sourcealignment}{Official Lab Manual --- Lab VIII}
\textbf{Objective.} Understand BSTs, heaps, AVL/height balancing, and B-trees and apply those structures to problems.

\textbf{Outcomes.} Learn BST, heap, and B-tree structures and implement/apply them computationally.

\textbf{Problem directions.} Create max/min heaps; insert/delete heap elements; solve BST problems; practice B-tree creation, insertion, and removal.
\end{sourcealignment}

\begin{labproject}{Compare ordering rules across specialized trees}
\textbf{Coverage.} Chapter 8.

\begin{lstlisting}
// BST insertion order
// 50, 30, 70, 20, 40, 60, 80
// inorder should be sorted.

// Max-heap insertion order
// 30, 10, 50, 20, 40
// print the heap array after every insertion.
\end{lstlisting}

\begin{tryit}
Insert 10,20,30,40,50 into a plain BST and draw the resulting height. Compare that shape with a balanced tree. Then create both a min-heap and max-heap from the same keys.
\end{tryit}

\begin{labmode}
The supplied course sources name B-tree creation/insertion/removal but do not provide the order/minimum-degree rules needed for a unique implementation. Follow the B-tree parameters specified by your instructor rather than inventing them.
\end{labmode}
\end{labproject}

\begin{verification}
\begin{itemize}
  \item BST inorder output is sorted for the chosen key policy.
  \item A max-heap parent is not smaller than either child; a min-heap parent is not larger than either child.
  \item Heap insertion/deletion restores the heap property.
  \item Complexity claims are related to tree height, not only node count.
\end{itemize}
\end{verification}

\begin{labdeliverable}
Submit a BST program, one heap implementation with insertion/deletion, and instructor-directed B-tree exercises where assigned.
\end{labdeliverable}
